The transport of electrons, phonon\cite{chen2018manipulation} and spins\cite{zutic2004spintronics,hirohata2020review} with/without the couplings amongst these three subsystems,  is a critical material property to be understood for developing new (nano)electronic technology through material innovations. The transport phenomena are associated with the dynamics of the carriers under external stimulus (\emph{e.g.} electric and/or magnetic field, temperature gradient) that induces an imbalance in the statistical distributions of the carrier across the sample.\cite{onsager1931reciprocal,onsager1931reciprocal2} It is therefore unsurprising that a plethora of theoretical approaches have been established to model the transport phenomena in different materials across a range of scales. 

In the context of atomistic modellings, for instances, in molecular and nano electronics with a device at the atomistic length scale, the quantum mechanical nature of the carrier plays an important role in determining the corresponding transport behaviours such that the transport phenomena is often interpreted as the time--evolution of the particle wavefunction.\cite{economou2006green} In this regime, the transport phenomena is often modeled with the Green's function approach under either the Landauer--B\"{u}ttiker\cite{buttiker1986four}  or the Kelydsh\cite{jauho1994time} formalism. The former is usually applied for modeling the phase coherent transport.\cite{xue2001charge,datta2005quantum} In contrast, the latter approach is more versatile to include the contributions from many--body  effects in modeling the inelastic transports.\cite{lu2007coupled,galperin2006resonant,liu2006inelastic} It has been combined with DFT for modelling the transport behaviors in realistic atomistic transport junctions,\cite{frederiksen2007inelastic,stradi2016general} but this method is highly specialised and its wide adaption, for example, in high--throughput computational screening is not common.

Charge transports in (crystalline) organic semiconductors, on the other hand, is most often interpreted as the hopping of electrons from one molecule to a neighboring one that are bounded together by the weak van der Waals interactions.\cite{coropceanu2007charge} The probability for the hopping to occur is characterised by the transfer integral $J_{ij}$ between the frontier orbitals of the two neighboring molecules, which itself is strongly correlated with the orbital overlaps $S_{ij}$.\cite{gajdos2014ultrafast} The latter provides a workable heuristic rule that is often applied, for instance, in guiding the design of organic semiconductors with high charge mobilities.\cite{li1998highly,takimiya2011thienoacene,gomez2017hydrogen,campbell2017predicted,yang2018large} Furthermore, the acceptance/donation of a charge carriers by a molecule will change its vibrational energy landscape $E(\bm{R})$ as a function of its atomic coordinates $\bm{R}$, resulted in the additional free energy cost associated with the charge hopping (most commonly referred to as the reorganisation energy $\lambda$). This `localised polaron' model is mathematically formulated in the Markus theory of charge transfer,\cite{marcus1993electron} in which the charge transfer rate between two molecule is given by
\begin{equation}\label{Markus}
\mathcal{W}_{ij}=\frac{|J_{ij}|^2}{\hbar}\sqrt{\frac{\pi}{\lambda k_B T}}\exp\left(-\frac{\lambda}{4k_{B}T}\right).
\end{equation}
The carrier mobility can then be obtained by using \cref{Markus} as inputs for  kinetic Monte Carlo simulations to model the diffusion pathways of charge carriers.\cite{coropceanu2007charge} Despite many known limitations,\cite{shuai2020applying} this model has widely be used in modelling the (upper bound limits of)  carrier mobilities in soft organic materials due to its conceptual simplicity. 

For inorganic solids, where the neighboring atoms are bounded by strong ionic/covalent interactions, their intrinsic  transport properties can be derived from their electronic band structures and/or densities--of--states,\cite{madsen2006boltztrap} which are readily accessible from most electronic structure calculations for periodic solids (for phonon transport, see discussions in  \cref{sec:anharmonic_vibs}). According to the Boltzmann transport theory, the electric current under an external gradient of temperature, electric and magnetic field is:
\begin{equation}\label{eq:boltzmanneq}
    j_i=\sigma_{ij}\bm{\mathcal{E}}_{j}+\sigma_{ijk}\bm{\mathcal{E}}_{j}\bm{\mathcal{B}}_{k}+v_{ij}\nabla_j T+\cdots
\end{equation}
From the topologies of the electronic bands, the group velocities and the tensor of inverse effective masses can be determined:
\begin{equation*}
    \bm{v}_\alpha(i,\bm{k})=\frac{1}{\hbar}\frac{\partial \varepsilon(i,\bm{k})}{\partial \bm{k}_{\alpha}}, \qquad \bm{M}^{-1}_{\alpha\beta}(i,\bm{k})=\frac{1}{\hbar^2}\frac{\partial^2 \varepsilon(i,\bm{k})}{\partial \bm{k}_{\alpha}\partial \bm{k}_\beta}.
\end{equation*}
Under the relaxation time approximation, the band--dependent conductivity tensors are then given by: 
\begin{equation}
\label{eq:e_cond_tensor}
    \bm{\sigma}_{\alpha\beta}(i,\bm{k})=e^2\tau_{i,\bm{k}}\bm{v}_\alpha(i,\bm{k})v_\beta(i,\bm{k}),
\end{equation}
\begin{equation}
\label{eq:hall_cond_tensor}
    \bm{\sigma}_{\alpha\beta\gamma}(i,\bm{k})=e^3\tau_{i,\bm{k}}^2\varepsilon_{\gamma uv}\bm{v}_\alpha(i,\bm{k})v_v(i,\bm{k})\bm{M}^{-1}_{\beta u},
\end{equation}
in which $\varepsilon_{ijk}$ is  the Levi--Civita symbol, and the band--dependent relaxation time $\tau_{i,\bm{k}}$ is usually assumed to be constant in most modelling cases. By further integrating the conductivity tensors over $\bm{k}$--points with a given energy, the density--of--states projected conductivity tensors $\bm{\sigma}_{\alpha\beta}(\varepsilon)$ and $\bm{\sigma}_{\alpha\beta\gamma}(\varepsilon)$ can be determined, from which the temperature and chemical potential dependent conductivity tensors [c.f. \cref{eq:boltzmanneq}] can be calculated as
\begin{subequations}
\begin{align}
&\mbox{electronic conductivity:}\nonumber  \\  &\sigma_{\alpha\beta}(T,\mu)=\frac{1}{\Omega}\int\bm{\sigma}_{\alpha\beta}(\varepsilon)\left[-\frac{\partial f_{\mu}(T,\varepsilon)}{\partial \varepsilon}\right]\mbox{d}\varepsilon, \label{electronic_cond}\\
&\mbox{electronic energy conductivity:}\nonumber\\     &\nu_{\alpha\beta}(T,\mu)=\frac{1}{eT\Omega}\int\bm{\sigma}_{\alpha\beta}(\varepsilon)(\varepsilon-\mu)\left[-\frac{\partial f_{\mu}(T,\varepsilon)}{\partial \varepsilon}\right]\mbox{d}\varepsilon, \label{electronic_energy_cond}\\
&\mbox{electronic thermal conductivity:}\nonumber\\    &\kappa_{\alpha\beta}^{e}(T,\mu)=\frac{1}{e^2T\Omega}\int\bm{\sigma}_{\alpha\beta}(\varepsilon)(\varepsilon-\mu)^2\left[-\frac{\partial f_{\mu}(T,\varepsilon)}{\partial \varepsilon}\right]\mbox{d}\varepsilon,\label{electronic_th_cond}\\
&\mbox{Hall conductivity:}\nonumber\\    &\sigma_{\alpha\beta\gamma}(T,\mu)=\frac{1}{\Omega}\int\bm{\sigma}_{\alpha\beta\gamma}(\varepsilon)\left[-\frac{\partial f_{\mu}(T,\varepsilon)}{\partial \varepsilon}\right]\mbox{d}\varepsilon, \label{Hall_cond}
\end{align}
\end{subequations}
where $f_{\mu}(T,\varepsilon)$ is the Fermi--Dirac distribution function. In particular, the total thermal conductivity of a given material can be obtained as a sum of its electronic thermal conductivity [\cref{electronic_th_cond}] and lattice thermal conductivity, providing a theoretical basis for the computational discoveries of thermoelectric materials.\cite{yang2008evaluation,hong2016full,zhang2017titanium,wang2017exceptional}

It can be seen that the above approach based on the Boltzmann transport equations establishes a nice link to bridge atomistic electronic structure simulation and the transport properties at the macroscopic limit. It avoids the complexity in considering the `split--over' of the electronic wavefunctions from the contact into the channel, which requires partitioning the systems into contact and channel regions.\cite{datta2005quantum} Although in its lowest level of approximations, detailed carrier scattering mechanisms (such as electron--phonon and impurity scatterings) are ignored, good agreements between the theoretically and experimentally determined transport characters had been demonstrated for materials such as  PbTe,\cite{singh2010doping} SnSe,\cite{zhao2016ultrahigh} Mg$_3$Sb$_{1.5}$Bi$_{0.5}$,\cite{zhang2017high} as well as a selected set of elemental and binary compound semiconductors.\cite{ricci2017ab}  Recent advances in material theory and availability in computational resources have seen more elaborated methods being developed and applied for determining the carrier relaxation times contributed by different scattering from first principles, enabling a more accurate prediction in the transport properties of exotic materials.\cite{krishnaswamy2017first,ganose2021efficient,jung2021prediction}

